\begin{exercise}[Derivative power counting for Nambu--Goldstone bosons]{I-CH09-023}
In four dimensions, normalize Nambu--Goldstone fields by
$U=\exp(\ii\pi^aT^a/f)$ and take the leading action to contain two
derivatives and a coefficient $f^2$.  Write the general symmetry-preserving
effective Lagrangian as a sum of vertices with $d_v=2,4,6,\ldots$ derivatives.

For a connected graph with $L$ loops and vertices $v$, derive Weinberg's
momentum-counting formula
\[
  \omega=2+2L+\sum_v(d_v-2).
\]
Show that an $E$-point amplitude has the form
\[
  \mathcal A_E\sim\frac{p^2}{f^{E-2}}
  \left(\frac{p}{4\pi f}\right)^{2L}
  \prod_v\left(\frac{p}{f}\right)^{d_v-2}
\]
up to dimensionless low-energy constants and group factors.  Conclude that
higher-derivative operators and loops give a systematic low-momentum
expansion.  Identify its range of validity, explain how analytic loop terms
renormalize higher operators, and check the Adler-zero and $p/f\to0$ limits.
\end{exercise}
